\begin{table}[H]
\centering
\caption{Relative performance comparison by MEAN_SPL using the Friedman test with post-hoc Nemenyi analysis.}
\label{tab:stats_mean_spl}
\begin{threeparttable}
\begin{tabular}{lccccc}
\toprule
Model & Mean MEAN_SPL & Std. & Avg. rank & Sig. wins & Sig. losses \\
\midrule
SetTransformer\_Q & 0.2109 & 0.0536 & 3.008 & 6 & 0 \\
M1\_AggHistFutureSummary & 0.2129 & 0.0567 & 3.074 & 6 & 0 \\
DeepSets\_Q & 0.2145 & 0.0563 & 3.314 & 6 & 0 \\
M0\_AggHistOnly & 0.2187 & 0.0596 & 3.727 & 6 & 0 \\
M3\_FullSkuTemporalCNN\_Q & 0.2183 & 0.0591 & 4.041 & 6 & 0 \\
BottomUpGlobalHistGB & 0.2585 & 0.1028 & 6.471 & 4 & 5 \\
AggregateHistGB & 0.2688 & 0.1058 & 7.587 & 2 & 5 \\
ChildSummaryHistGB & 0.2724 & 0.1076 & 7.909 & 2 & 6 \\
AggregateElasticNet & 0.3227 & 0.1743 & 8.025 & 1 & 6 \\
ChildSummaryElasticNet & 0.3469 & 0.1950 & 9.231 & 0 & 8 \\
SeasonalNaive & 0.3356 & 0.1450 & 9.612 & 0 & 9 \\
\bottomrule
\end{tabular}
\begin{tablenotes}
\footnotesize \item The Friedman test uses 121 aligned series-level blocks (task,agg\_id). The test statistic is $\chi^2=744.6491$ with $p=1.62e-153$. Average ranks are computed with lower MEAN_SPL indicating better performance. Nemenyi significant wins/losses are counted at $\alpha=0.10$ using critical difference $CD=1.2697$.
\end{tablenotes}
\end{threeparttable}
\end{table}
